% =====================================================================
% NSF CDS&E Proposal (PD 24-8084) — standard NSF PAPPG 24-1 format.
% CDS&E is a META-PROGRAM: submit into a participating division program.
% After the CISE drop (PD 24-8084), CDS&E routes through ENG (CBET/CMMI)
% + MPS only; ENG/CBET (biomedical / data-enabled discovery) is the most
% plausible home for this topic. Title MUST begin "CDS&E:".
% >>> Confirm program fit with a 1-paragraph concept note to the PO BEFORE
%     submitting; the Project Description is portable to SaTC / NLM. <<<
%
% Formatting (PAPPG 24-1, current June 2026):
%   - >= 11pt Times/Computer Modern; >= 1in margins; <= 6 lines/vertical inch
%   - Project Summary 1 pg; Project Description 15 pg; References no limit
%   - Biosketch + Current & Pending: SciENcv only (NOT LaTeX)
%
% Compile: pdflatex main; bibtex main; pdflatex main; pdflatex main
%   (references live in references.bib)
% =====================================================================
\documentclass[11pt]{article}

\usepackage[margin=1in, letterpaper]{geometry}
\usepackage[utf8]{inputenc}
\usepackage[T1]{fontenc}
\usepackage{graphicx}
\usepackage{booktabs}
\usepackage{enumitem}
\setlist{nosep, leftmargin=1.4em}
\usepackage[hidelinks]{hyperref}
\usepackage{titlesec}
\titlespacing*{\section}{0pt}{1.2ex plus .4ex}{0.6ex}
\titlespacing*{\subsection}{0pt}{1ex plus .3ex}{0.4ex}
\titleformat{\section}{\normalfont\large\bfseries}{\thesection}{0.6em}{}
\titleformat{\subsection}{\normalfont\normalsize\bfseries}{\thesubsection}{0.6em}{}

\newcommand{\ptitle}{CDS\&E: The Cost of Patient Autonomy in Cross-Silo
Federated Learning --- Measuring the Computational Price of Dynamic,
Revocable Consent in Federated Clinical Prediction}

\begin{document}

% =====================================================================
% PROJECT SUMMARY — 1 page. Third person. Three labeled parts.
% =====================================================================
\begin{center}
  \textbf{Project Summary}\\[2pt]\textbf{\ptitle}
\end{center}

\noindent\textbf{Overview.}
Federated learning (FL) lets hospitals train a shared model without moving
patient data, and a growing literature adds a blockchain layer so that patients
can grant and revoke consent to participate. The field treats that capability as
an unqualified good. What it has \emph{not} done is measure its price. Giving
each patient fine-grained, dynamic, revocable control over their participation in
a federated model is not free: it imposes on-chain transaction and bandwidth
overhead on every training round, and---because consented data becomes a moving
target---it perturbs the very statistical population the model learns from,
costing model utility. This project designs, builds, and rigorously measures a
prototype \emph{consent-gated cross-silo federated learning} system in order to
quantify, for the first time, the two-axis \emph{cost of patient autonomy}: the
\emph{systems cost} (governance overhead per round) and the \emph{utility cost}
(accuracy degradation as consent churns). The headline deliverable is not a more
accurate model; it is a set of empirically grounded \emph{cost--utility tradeoff
curves} and the configurations and breaking points at which patient-controlled
federated clinical prediction remains practical. The system uses an on-chain
two-tier
consent layer (silo admission $+$ per-patient consent) over an encrypted
off-chain store, drives a federated 30-day hospital-readmission classifier over
a real cohort of roughly 102{,}000 inpatient encounters (UCI Diabetes
130-US-hospitals), and is exercised under controlled consent-churn regimes, with
synthetic FHIR patient populations scaling the systems axis, where signal
quality is irrelevant. The team has independently prototyped and
published each constituent component, and has already executed the utility-axis
measurement---confirming the central churn-regime hypothesis and, via a
model-capacity replication, its independence of the client model---giving
strong feasibility evidence for the proposed integration and measurement.

\medskip
\noindent\textbf{Intellectual Merit.}
The project advances computational and data-enabled methods for trustworthy
multi-institutional learning by reframing patient-consent governance from a
feature to be \emph{built} into a cost to be \emph{measured}. Its contributions
are: (i) the first dual-axis \emph{characterization} of consent-governance
overhead---decomposing each federated round into compute versus on-chain
governance cost, and accuracy loss as a function of consent dynamics---where the
prior literature reports the architecture but not its price; (ii) a separation of
three physically distinct ``churn'' regimes that the literature conflates
(transient/rotating departure, permanent withdrawal, and whole-silo exit) and
the finding of where each does and does not degrade a federated model; and (iii)
a generalizable, FHIR-aligned, multi-project consent framework (a
patient~$\times$~study matrix) that demonstrates governance over patients
enrolled in several federated studies at once. The work is explicitly a
measurement and systems contribution, not a new learning algorithm, and is
designed to surface \emph{non-obvious} findings---e.g., overhead thresholds
beyond which per-round on-chain consent is infeasible, and the result that
whole-silo departure damages utility far more than individual patient churn.
Executed five-seed experiments on the primary cohort already establish the
latter: a positive-heavy whole-silo exit costs $-0.059$ AUROC against $-0.039$
for a count-matched random control---isolating the damage to \emph{who} leaves
rather than \emph{how many}---and a replication with a higher-capacity neural
client reproduces every delta within noise, confirming the effect is a property
of the problem rather than of the model.

\medskip
\noindent\textbf{Broader Impacts.}
Unplanned hospital readmissions cost the U.S.\ health system billions of dollars
annually and fall hardest on underserved populations; readmission-risk
prediction is a standard hospital analytics task whose accuracy depends on
diverse, multi-institutional data that privacy regulation keeps siloed. The
project's Year-3 extension targets cervical-cancer screening---the fourth most
common cancer in women worldwide, with the overwhelming burden in low- and
middle-income regions where screening infrastructure is weakest. In both
settings, knowing the real computational cost of
patient-controlled, privacy-respecting collaborative learning is a precondition
for deploying it where it would matter most: a system whose governance overhead
is unmeasured cannot be sized, budgeted, or trusted. By delivering tradeoff
curves and breaking points rather than a single point estimate, the project gives
system designers, IRBs, and health-data regulators an evidence base for setting
consent-refresh policies. The project trains graduate and undergraduate
researchers at a primarily-undergraduate-serving institution in the rare
intersection of distributed ML, data engineering, and applied cryptography, and
releases an open, reproducible reference implementation and synthetic benchmark.

\newpage

% =====================================================================
% PROJECT DESCRIPTION — 15 pages max.
% =====================================================================
\begin{center}\textbf{Project Description}\\[2pt]\textbf{\ptitle}\end{center}

\section{Introduction and Motivation}
Clinical prediction models---which patients will be readmitted within 30 days of
discharge, which screening results warrant follow-up---can extend the capacity
of strained health systems, but accurate models
need large, diverse training data, and that data is locked inside individual
hospitals: patient-privacy regulation, institutional liability, and the
impracticality of moving medical records across organizations forbid central
pooling. The stakes span the care continuum, from unplanned readmissions that
cost the U.S.\ system billions annually to cervical-cancer screening, whose
burden falls overwhelmingly on low- and middle-income regions where screening
infrastructure is weakest~\cite{globocan}. Federated learning (FL)~\cite{fedavg,kairouz} is the standard response:
keep data local, share only model updates. FL has been demonstrated for
multi-institutional medicine without sharing patient data~\cite{sheller}, and a
rapidly growing line of work adds a \emph{blockchain} layer so that
participation---by institutions and increasingly by individual patients---can be
granted, revoked, and audited on an immutable ledger~\cite{mdpi2025,oxford2025}.

\paragraph{The gap: a capability built but never priced.}
The consent-governance literature has converged on a clear architectural answer
and a clear normative claim---patients \emph{should} be able to control, in real
time, whether their data trains a given model---and recent systems implement
exactly that~\cite{oxford2025,mdpi2025}. What the literature has \emph{not} done
is measure what that control costs. Patient autonomy in a federated system is not
free along two independent axes:
\begin{itemize}
\item \textbf{Systems cost.} Every consent grant or revocation is an on-chain
transaction with latency and (on public chains) monetary cost; every training
round must resolve the current eligible-patient set before it can proceed. On a
public chain, a single permission transaction can take seconds to minutes under
congestion---our own prior measurements range from $\sim$9\,s to $\sim$108\,s on
a public testnet (\S\ref{sec:feas}). If consent must be re-checked each round,
governance can dominate wall-clock time.
\item \textbf{Utility cost.} Consent that can be revoked at any time means the
training population is \emph{non-stationary}. Patients (and entire silos) entering
and leaving between rounds perturb the data distribution---precisely the non-IID
condition under which FedAvg is known to lose accuracy~\cite{zhao,fedprox}. The
right to leave therefore has a measurable price in model quality.
\end{itemize}
These two costs are in tension with each other and with the autonomy they buy.
Cheaper governance (amortizing consent checks across many rounds) means consent
is honored more slowly; stricter, per-round consent enforcement means higher
overhead and a more volatile training population. \emph{Nobody has mapped this
frontier.}

\paragraph{Thesis and research question.}
This project's contribution is a \emph{measurement}, not a new architecture or a
new learning algorithm. We ask:
\begin{quote}
\emph{What does it cost---in on-chain and bandwidth overhead, and in model
utility---to give patients fine-grained, dynamic, revocable consent in a
cross-silo federated clinical-prediction system, and at what configurations and
breaking points does patient-controlled federated learning remain practical?}
\end{quote}
The headline deliverable is a family of \emph{cost--utility tradeoff curves} and
the breaking points on them, not a peak-accuracy number. Chasing accuracy would
optimize the wrong thing; the scientific weight of the work rests on
characterizing the \emph{deltas}---how cost and utility move as the parameters of
patient autonomy (consent-churn rate, churn regime, consent-refresh cadence,
chain backend, number of concurrent studies) are swept. We design the study
specifically to surface \emph{non-obvious} findings (\S\ref{sec:rq}): for
instance, an overhead threshold beyond which per-round on-chain consent is
infeasible and must be amortized; the result that transient/rotating churn may
\emph{not} degrade accuracy while whole-silo departure sharply does; and the cost
profile of honoring a revocation that demands the model \emph{forget} a
withdrawn patient.

\paragraph{Why now, and why this team.}
The architectural pieces---blockchain-mediated consent, encrypted off-chain
storage, cross-silo FL, scalable medical ML---now individually exist, several of
them in this team's own published work (\S\ref{sec:feas}). The field has built
the machine; it has not put it on a dynamometer. This proposal supplies the
measurement methodology and the instrument.

\section{Research Questions and Hypotheses}\label{sec:rq}
We organize the work around four research questions. Each is paired with a
falsifiable hypothesis and an explicit ``non-obvious finding'' target, because an
integration of known parts earns its keep only through what the evaluation
reveals---``more dropout lowers accuracy'' is not a finding.

\begin{description}
\item[RQ1 (Systems cost).] How does on-chain consent governance overhead scale
with consent-churn rate, consent-refresh cadence, federation size, and chain
backend (local / permissioned / public)? \emph{Hypothesis H1:} per-round
governance cost is dominated by consent-set resolution and grows with churn rate,
crossing a backend-dependent threshold beyond which per-round enforcement is
infeasible and consent must be amortized across rounds. \emph{Non-obvious target:}
locate that threshold and show the amortization--staleness tradeoff it forces.

\item[RQ2 (Utility cost \& churn regimes).] How does final model utility degrade
as a function of consent dynamics, and do physically distinct churn regimes
degrade it differently? We separate three regimes the literature conflates: (a)
\emph{transient/rotating} churn (patients leave and rejoin), (b) \emph{permanent
withdrawal} (further split into random and outcome-biased variants), (c)
\emph{whole-silo departure}. \emph{Hypothesis H2:} the three are
not interchangeable---transient churn may be near-harmless because contributions
re-enter, whereas whole-silo departure inflicts a non-IID shock that dominates
all per-patient effects. \emph{Non-obvious target:} show that ``who leaves''
matters more than ``how many leave,'' overturning the intuitive headcount model.
\emph{Status: confirmed---a five-seed study on the primary cohort separates the
regimes and passes the isolation test, and an executed model-capacity
replication (MLP client) reproduces the deltas within noise
(\S\ref{sec:feas}).}

\item[RQ3 (The frontier).] Where is the joint cost--utility frontier, and what
configurations keep the system in a practical operating region? \emph{Hypothesis
H3:} there exist consent-refresh cadences that recover most of the utility of
per-round enforcement at a fraction of the governance cost, i.e.\ the frontier is
sharply ``knee-shaped'' rather than linear. \emph{Non-obvious target:} identify
the knee and the deployment regimes it implies.

\item[RQ4 (Multi-project governance).] Can a single consent layer govern patients
enrolled in \emph{several} federated studies at once (a patient~$\times$~study
matrix), and what is the marginal governance cost per additional study?
\emph{Hypothesis H4:} multi-project consent adds bounded, sub-linear on-chain
overhead because consent state is shared infrastructure, demonstrating that the
framework generalizes beyond a single disease model without building a second
one.
\end{description}

\section{Results from Prior NSF Support}
The PI has no prior NSF support. (Senior personnel to confirm; if any held NSF
awards in the past five years, summarize here per PAPPG 24-1 with award number,
amount, period, and Intellectual Merit / Broader Impacts outcomes.)

\section{Preliminary Studies and Feasibility}\label{sec:feas}
Each layer of the proposed system has already been prototyped and published by
members of the team. This is the principal feasibility evidence: the project is an
\emph{integration and measurement} of validated components, not a from-scratch
build, which is what makes the ambitious measurement program achievable in scope.

\subsection{Scalable cervical-cancer classification (the compute-placement rule and Year-3 image baseline)}
The team built a scalable cervical-cancer detection pipeline on Apache
Spark/SparkML and benchmarked it against scikit-learn and a convolutional
baseline across two modalities~\cite{cervical}. On the tabular cervical-cancer
risk dataset, a Spark Random Forest reached 0.96 accuracy (0.99 precision, 0.97
recall) after ADASYN balancing of a strongly imbalanced cohort; in a scalability
study (data scaled from 0.27 to 217\,MB), Spark crossed over to roughly
$3.3\times$ faster than scikit-learn as data volume grew. For Pap-smear images
(SipakMed), a CNN reached 92.75\% where Spark's general classifiers plateaued
near 73\%. \emph{This supplies---critically---the
empirically grounded design rule that distributed/Spark compute earns its place
on large image data but is overkill on small tabular data, plus the Year-3 image
task and baseline. We honor that rule
below by using plain tabular federated training in Year~1 and deferring
Spark/images to Year~3.}

\subsection{Patient-controlled health records on blockchain (the consent layer)}
The team built a patient-centric electronic health record system in which a
Python application mediates an Ethereum smart contract and IPFS
storage~\cite{ehr}. Records are encrypted (Fernet) and stored on IPFS; the
contract holds the content identifier, owner wallet, authorized-viewer list, and
key, and exposes \texttt{grantPermission}, \texttt{rescindPermission},
\texttt{viewInformation}, and \texttt{isViewerAuthorized}. Measured performance:
IPFS retrieval typically under 0.3\,s; on a local chain (Ganache), add-permission
averaged 0.052\,s / 0.00099\,ETH and revoke 0.071\,s / 0.00119\,ETH; on a public
testnet (Sepolia), permission transactions ranged from $\sim$9\,s to
$\sim$108\,s under congestion. \emph{This supplies the consent-layer design, the
grant/revoke semantics, and---the key instrument for this project---a validated
methodology for measuring on-chain governance cost (latency and gas per
operation), which is exactly what RQ1 requires.}

\subsection{Scalable predictive modeling of multi-institutional data \& honest benchmarking}
In a third line of work the team modeled lead time across a large
multi-institutional healthcare supply chain with Random Forest regression
($R^2 \ge 0.87$) and showed feature importance shifting between normal and crisis
regimes~\cite{nhan2024optimizing}---a documented multi-institution regime-shift, the
empirical cousin of the non-IID shock studied in RQ2. A separate unified benchmark
contributed the team's methodology for honest model comparison---calibration
(expected calibration error), McNemar significance testing, and the finding that
well-chosen features can outweigh large increases in model
capacity~\cite{tennis}. \emph{These supply the evaluation rigor (calibration,
significance testing) that keeps the tradeoff curves defensible rather than
anecdotal.}

\subsection{Executed preliminary study: the utility axis on real data}
\label{sec:prelim}
Beyond component feasibility, the utility-cost axis (RQ2) has already been
executed in preliminary form on the primary cohort (UCI Diabetes
130-US-hospitals: 101{,}766 real inpatient encounters, 30-day readmission, 11.2\%
positive; \S\ref{sec:data}). A multi-model centralized benchmark (logistic
regression, random forest, gradient boosting, XGBoost, LightGBM; five-fold CV)
establishes the ceiling at AUROC $\approx 0.677$---inside the published
0.667--0.70 band for this task---with logistic regression, the best
weight-averageable model, at $0.666$ setting the federated ceiling. A five-method
federated comparison at $K{=}3$ (FedAvg, FedProx, FedAvgM, FedAdam, SCAFFOLD)
shows aggregation choice is immaterial at IID and moderate skew (AUROC
$0.666$--$0.670$), with SCAFFOLD collapsing under severe heterogeneity
(Dirichlet $\alpha{=}0.1$: $0.625$)---empirically justifying class-weighted
FedAvg as the instrument. The three-regime churn study (five seeds, churn to
70\%) then yields, as $\Delta$AUROC against the no-churn baseline: transient
churn $-0.003$; permanent-random $-0.007$; permanent-biased (positive-outcome
patients leave first) $-0.021$; whole-silo exit of a positive-light silo
$-0.010$; whole-silo exit of the positive-heavy silo $-0.059$. The decisive
comparison is a \textbf{count-matched isolation test}: removing the same number
of patients as the positive-heavy silo exit, but drawn uniformly at random,
costs only $-0.039$---so identity, not headcount, carries the damage,
\emph{confirming H2}. A subsequent \textbf{model-capacity replication} reran
the entire churn grid with two clients on identical partitions and schedules
(200 runs): the linear client, and a one-hidden-layer MLP (64 units) that
centralizes at AUROC $0.670$, above the linear ceiling of $0.666$. Every delta
matches within one standard deviation (positive-heavy silo exit $-0.064$ on the
MLP vs.\ $-0.059$; count-matched control $-0.051$ vs.\ $-0.039$), so the
churn-cost structure is model-independent and H2 survives the model swap.
These results de-risk the utility axis entirely; the funded work scales it
(finer rate sweeps, calibration and significance testing per~\cite{tennis}),
couples it to the systems axis, and assembles the frontier.

\section{Background and Related Work}
\paragraph{Federated learning and its non-IID weakness.}
FedAvg~\cite{fedavg} is the canonical cross-silo algorithm and our Year-1
baseline. Its sensitivity to non-identically distributed (non-IID) client data is
well documented~\cite{zhao}; FedProx~\cite{fedprox} and
SCAFFOLD~\cite{scaffold} were designed to mitigate it. We treat FedAvg's non-IID
sensitivity not as a flaw to engineer away but as the \emph{measurement
instrument}: it is precisely what makes consent churn visible as a utility cost.
Our executed five-method comparison (\S\ref{sec:prelim}) already shows the
aggregator choice is immaterial at $K{=}3$ short of pathological heterogeneity,
and the model-capacity check has likewise been executed---an MLP averageable
client reproduces the degradation magnitudes---so the Year-2 robustness check
shifts to replication on real hospital silos (eICU).

\paragraph{Blockchain $+$ FL in healthcare.}
This is now a heavily reviewed area, but across it the blockchain's role is
overwhelmingly to \emph{verify or aggregate} model updates, provide
\emph{provenance}, or \emph{incentivize} participation---not to govern
fine-grained patient consent, and almost never accompanied by a measurement of
its cost~\cite{sheller}.

\paragraph{Patient-consent-gated FL (closest prior work).}
Two works are nearest. Kostick-Quenet et al.~\cite{oxford2025} argue, as legal
architecture, for per-patient consent smart contracts as gatekeepers of FL, but
provide \emph{no implementation and no measurement}. The MDPI Electronics
system~\cite{mdpi2025} \emph{implements} a blockchain-$+$-FL EHR-sharing
framework with dynamic differential privacy and federated distillation, and gates
participation---but it reports the \emph{architecture and security properties},
not a characterization of governance overhead or of how consent dynamics degrade
utility. \textbf{This is the precise opening:} the closest implemented system
builds the consent-gated machine but does not price it; the closest consent
framework is conceptual only. Our contribution is the measurement neither
provides---the cost of autonomy, on both axes, with breaking points.

\paragraph{FL for cervical cancer.}
A CNN-based federated system on SipakMed reports $\sim$94.4\% IID and $\sim$78.4\%
non-IID accuracy~\cite{joynab2024}---a 16-point non-IID gap that concretely sizes
the utility axis we measure, and a baseline for the Year-3 image work. We are
aware of no federated cervical system that is blockchain-consent-gated, and none
that measures the cost of that gating.

\paragraph{Synthetic data for the systems axis.}
Synthea~\cite{synthea} generates unlimited, FHIR-native synthetic patients. We
use it deliberately and narrowly (\S\ref{sec:data}): on the systems axis, signal
quality is irrelevant---only transaction volume and payload size matter---so
Synthea scales consent-governance traffic to hospital scale to locate the H1
infeasibility threshold. The utility axis, by contrast, runs entirely on real
encounter data, so degradation curves measure genuine clinical signal rather
than generator artifacts.

\paragraph{Privacy, unlearning, and where they sit.}
We are careful with the word ``privacy.'' FedAvg alone yields data-\emph{locality},
which is leaky, not privacy. Technical privacy means differential
privacy~\cite{dpsgd} (noise on gradients) or secure aggregation~\cite{secagg}
(cryptographic), composable layers \emph{on top} of FL. In Year~1 we therefore
claim ``consent governance and data autonomy,'' not privacy; DP/secure-aggregation
is a Year-2 thrust that earns the term and adds another knob (noise~$\rightarrow$
accuracy loss) to the same cost frontier. Honoring a permanent withdrawal can, in
the limit, require the model to \emph{forget} a contribution---machine
unlearning~\cite{unlearning}---which we scope as a measured stretch goal, not a
core deliverable.

\section{Research Plan}
The plan is organized as five thrusts mapped onto a three-year arc
(\S\ref{sec:timeline}). Thrusts~1--3 are the dissertation/Year-1 \emph{spine} and
answer RQ1--RQ4 on tabular data at statistically powered scale; Thrusts~4--5
deepen the frontier (privacy, robustness, images) and are the grant's Year-2/3
superset.

\subsection{Thrust 1 --- Consent-gated cross-silo FL testbed (the instrument)}
Build the measured system end to end on tabular data: (a) a FedAvg coordinator
over $K$ hospital silos training the readmission classifier---the churn-aware
trainer and consent-schedule library are already built and validated
(\S\ref{sec:prelim}); (b) a
\emph{two-tier} on-chain consent layer---a silo-registry contract (admission,
mostly static) and a patient-consent contract (dynamic, grant/revoke), extending
the team's prior contract~\cite{ehr}; (c) encrypted off-chain storage (IPFS) with
the chain holding only \texttt{\{patientID, consentStatus, CID, key,
authorizedViewers\}}---\emph{never} clinical data on-chain, which would be a
privacy disaster given ledger immutability and replication; and (d) the
coordinator query path that, before each round (or each refresh interval),
resolves the eligible silo set and per-silo consented record set. \emph{Deliverable:}
a correct, auditable, fully instrumented testbed---the dynamometer for all
measurement below. \emph{Risk/mitigation:} public-chain latency makes per-round
enforcement impractical; we instrument local, permissioned, and public backends
and treat the consent-refresh cadence as a swept parameter rather than a fixed
choice.

\subsection{Thrust 2 --- The systems-cost axis (RQ1)}
Instrument every federated round to decompose wall-clock and monetary cost into
\emph{compute} versus \emph{governance} (consent-set resolution, on-chain
grant/revoke, content retrieval), reusing the team's gas/latency methodology~\cite{ehr}.
Sweep churn rate, refresh cadence, federation size $K$, and chain backend.
\emph{Deliverable:} governance-overhead curves and the H1 infeasibility
threshold; the amortization--staleness tradeoff. \emph{Measurement integrity:} we
report per-operation distributions (not just means) because public-chain latency
is heavy-tailed, and we log any capping or sampling explicitly.

\subsection{Thrust 3 --- The utility-cost axis: the churn study (RQ2)}
Drive the testbed under controlled consent dynamics and measure accuracy
degradation. We sweep churn \emph{rate} (10/30/50/70\%) and, crucially, run three
\emph{distinct} regimes rather than conflating them:
\begin{enumerate}
\item \textbf{Transient/rotating churn} --- patients withdraw consent and later
restore it; contributions re-enter the population. (Preliminary: near-harmless,
$-0.003$ AUROC at 70\%.)
\item \textbf{Permanent withdrawal} --- contributions leave for good, split into
random and outcome-biased variants. (Preliminary: random is cheap, $-0.007$;
biased withdrawal costs $3\times$ more, $-0.021$---distribution shift, not
headcount, drives the price. The natural home for the unlearning stretch goal.)
\item \textbf{Whole-silo departure} --- an institution exits, removing a whole
slice of the distribution at once. (Preliminary: the dominant regime, $-0.059$
for a positive-heavy silo under severe skew.)
\end{enumerate}
\emph{Deliverable:} the full powered degradation curves per regime with
calibration and significance testing, extending the executed study
(\S\ref{sec:prelim}) that already confirms the H2 ``who leaves $>$ how many
leave'' finding via its count-matched isolation test and its model-independence
via the MLP replication; funded work adds finer rate sweeps, calibration and
significance testing, and real-hospital silo replication (\S\ref{sec:data}).

\subsection{Thrust 4 --- Multi-project consent \& the frontier (RQ3, RQ4)}
Extend the consent layer to a patient~$\times$~study matrix at the
contract/ledger level---a patient may be enrolled in, and independently withdraw
from, several federated studies---\emph{without} building a second disease model
(we instantiate additional ``studies'' as parameterized federated runs over the
same infrastructure). Measure marginal governance cost per added study (RQ4) and
assemble the joint cost--utility frontier, locating the H3 ``knee'' where a
relaxed refresh cadence recovers most utility at a fraction of the cost.
\emph{Deliverable:} the frontier and its practical operating region---the
proposal's central scientific artifact.

\subsection{Thrust 5 --- Earning ``privacy,'' robustness, and images (Year 2/3 superset)}
\begin{itemize}
\item \textbf{Privacy thrust (Year 2).} Add DP-SGD~\cite{dpsgd} and/or secure
aggregation~\cite{secagg} as composable layers and measure the noise~$\rightarrow$
accuracy cost as a third axis on the same frontier; this is the thrust that earns
the word ``privacy.''
\item \textbf{Robustness (Year 2).} Two robustness dimensions are already
closed---the aggregator (five methods compared, FedAvg justified) and model
capacity (an executed MLP replication reproduces every churn delta within
noise, \S\ref{sec:prelim})---so robustness shifts to replication on real
hospital silos (eICU), testing whether Thrust-3 degradation magnitudes are
artifacts of Dirichlet partitioning.
\item \textbf{Unlearning (stretch).} On permanent withdrawal, trigger federated
unlearning~\cite{unlearning} of the revoked contribution and measure its
cost---an open problem flagged but unaddressed by the prior consent-gating
work~\cite{oxford2025}.
\item \textbf{Images (Year 3).} Move to cervical images (SipakMed) where the heavy
multimodal/Spark pipeline~\cite{cervical} is \emph{honestly} justified, re-running
the cost--utility study at image scale against the federated-cervical
baseline~\cite{joynab2024}.
\end{itemize}

\section{Data and Materials}\label{sec:data}
A single powered primary dataset carries the utility axis; synthetic data is
confined to the systems axis, where signal quality is irrelevant by
construction:
\begin{itemize}
\item \textbf{UCI Diabetes 130-US-hospitals (primary; real signal, powered
$N$).} 101{,}766 real inpatient encounters from 130 U.S.\ hospitals
(1999--2008); binary 30-day-readmission label, 11.2\% positive ($\approx$11{,}000
positive cases). Large enough that per-silo churn deltas remain statistically
powered at 70\% churn, and real enough that degradation curves measure genuine
clinical signal. This dataset carries the entire utility axis (Thrust~3), as
the executed preliminary study demonstrates (\S\ref{sec:prelim}).
\item \textbf{Synthea~\cite{synthea}} (FHIR-native, unlimited $N$, weak ML
signal) carries the \emph{systems/cost} metrics only---signal quality is
irrelevant to bandwidth/overhead---scaling consent-governance traffic to
hospital scale for the H1 threshold. For FHIR oncology conformance in later
years we use mCODE-aligned profiles~\cite{mcode}.
\item \textbf{Real cervical risk data} (public 858-row clinical risk-factor
dataset, real signal) serves as a small \emph{signal-validity anchor} and
continuity with the team's published classifier~\cite{cervical}. At $\sim$286
records/silo across three silos (54 positive cases in total) it cannot power a
churn study---the quantitative fact that motivated the primary-dataset choice.
\item \textbf{Year 2/3: real hospital silos (eICU)} for replication with natural
institutional boundaries in place of Dirichlet partitions; \textbf{Year 3:
cervical images (SipakMed)}---large, real signal---justify
the image pipeline and the Spark/CNN stack honestly.
\end{itemize}
We use a \textbf{FHIR-aligned subset} (Patient, Observation, Condition; Media for
images) with LOINC/SNOMED coding---``aligned,'' not fully ``conformant'';
conformance is a Year-2/3 garnish, not a Year-1 gate. No real patient-identifying
data is used; synthetic consent populations drive the governance layer.

\section{Evaluation Plan and Metrics}
\textbf{Systems (cost axis):} per-round and per-operation on-chain latency and gas
on local vs.\ permissioned vs.\ public backends; bandwidth per round; end-to-end
round time decomposed into compute vs.\ governance; reported as distributions.
\textbf{Model (utility axis):} accuracy, precision/recall, F1, ROC-AUC, PR-AUC,
and
\emph{calibration} (expected calibration error), with McNemar tests for
significance vs.\ centralized and no-churn baselines~\cite{tennis}; explicit IID
vs.\ non-IID partitions. \textbf{Functional:} grant/deny matches policy at both
tiers; revocation propagates by the next eligible round. \textbf{Headline
artifacts:} the cost--utility tradeoff curves (per churn regime, per backend, per
refresh cadence), the RQ1 infeasibility threshold, and the RQ3 frontier knee.
\textbf{Reproducibility:} fixed seeds, released code/contracts, and the synthetic
benchmark.

\section{Timeline and Management}\label{sec:timeline}
\textbf{Year 1 (spine):} Thrust 1 testbed; Thrust 2 at Synthea scale; Thrust 3
completed on the primary cohort (already executed, including the
model-capacity replication, \S\ref{sec:prelim}); first cost--utility curves and
the three-regime churn finding---this is the
PhD dissertation's core. \textbf{Year 2:} Thrust 4 multi-project consent and the
frontier; Thrust 5 privacy (DP/secure-agg) and robustness (eICU real-silo
replication).
\textbf{Year 3:} Thrust 5 images (SipakMed), unlearning stretch goal, mCODE
conformance, open release and dissemination. The PI coordinates; a graduate
research assistant leads implementation; undergraduates contribute to data and
measurement pipelines (see Broader Impacts). A Gantt chart accompanies the
submission.

\section{Broader Impacts}
The deliverable---tradeoff curves and breaking points, not a single
accuracy number---gives system designers, IRBs, and health-data regulators the
evidence base to set consent-refresh and governance policy for
patient-controlled federated learning. That evidence base matters wherever
siloed data limits clinical prediction: for readmission risk, whose burden falls
hardest on underserved populations, and---in the project's Year-3
extension---for cervical screening, where systems must eventually be deployed in
the low- and middle-income regions bearing the highest burden with the weakest
infrastructure~\cite{globocan}. The
project provides research training at a primarily-undergraduate institution in
distributed ML, data engineering, and applied cryptography, deliberately involving
undergraduates in measurement and data pipelines. All software, smart contracts,
and the synthetic benchmark are released open-source with documentation for
reproducibility and reuse.

\section{Intellectual Merit (summary)}
The project delivers the first dual-axis \emph{measurement} of the cost of patient
autonomy in cross-silo federated learning: governance overhead and utility
degradation under dynamic, revocable consent, with the configurations and breaking
points where patient-controlled federated clinical prediction stays practical. It
separates
three conflated churn regimes---a separation its executed preliminary study
already demonstrates on 101{,}766 real encounters, confirming that \emph{who}
leaves matters more than \emph{how many}---generalizes via a multi-project
consent matrix, and
is grounded in the team's prior, published components---a data-enabled systems
contribution that prices a capability the field has built but never measured.

% =====================================================================
\newpage
\begin{center}\textbf{References Cited}\end{center}
\renewcommand{\refname}{}
\vspace{-2em}
{\small
\bibliographystyle{ieeetr}
\bibliography{references}
}

\end{document}
